% GENERATED FILE. Edit canonical lessons, then run npm run generate:books.
% canonical-source-hash: fnv1a64:c13d0640bcac9098
% canonical-lessons: BN-C100-sahos, BN-C100-protibha, BN-C100-jontro, BN-C100-soronjam, BN-C100-poddhoti

\chapter{Courage, Talent, Machine, Equipment, Method}
\label{ch:bn-courage-talent-machine-equipment-method}
\hlchaptermodality{100}

\begin{chapteropening}
\textbf{By the end of this chapter:} \emph{I can say courage, talent, a machine, equipment and a method.}

Complete the last lesson of chapter 100: i can say courage, talent, a machine, equipment and a method.
\end{chapteropening}

\section[\texorpdfstring{\bn{সাহস} (sāhôs)}{সাহস (sāhôs)}]{\bn{সাহস} (sāhôs) — courage}

\label{lesson:BN-C100-sahos}



\begin{quote}
Before the new one: say the Bengali for intelligence, then the Bengali for knowledge.
\end{quote}

\subsection*{You'll want to know: \bn{সাহস}}
\textbf{\bn{সাহস}} — \emph{sāhôs} — \textquotedblleft{}courage\textquotedblright{}.

\begin{grammarlens}[title={What you've built}]
One more word for this chapter.
\end{grammarlens}

\subsection*{Your turn}
\begin{itemize}
\raggedright
  \item \emph{Say it:} \emph{sāhôs}
  \item \emph{Say it:} \emph{sāhôs}, once more
  \item \emph{Say it:} say \emph{sāhôs}
  \item \emph{Recall:} say the Bengali for intelligence, then the Bengali for knowledge, then say \emph{sāhôs} again
\end{itemize}

\begin{tcolorbox}[breakable,colback=teal!4,colframe=teal!35!black,title={Before you move on}]
What does \bn{সাহস} mean? (\textquotedblleft{}Courage\textquotedblright{}.) Say it once more.
\end{tcolorbox}
\section[\texorpdfstring{\bn{প্রতিভা} (prôtibhā)}{প্রতিভা (prôtibhā)}]{\bn{প্রতিভা} (prôtibhā) — talent}

\label{lesson:BN-C100-protibha}



\begin{quote}
Before the new one: say the Bengali for knowledge, then the Bengali for courage.
\end{quote}

\subsection*{You'll want to know: \bn{প্রতিভা}}
\textbf{\bn{প্রতিভা}} — \emph{prôtibhā} — \textquotedblleft{}talent\textquotedblright{}.

\begin{grammarlens}[title={What you've built}]
One more word for this chapter.
\end{grammarlens}

\subsection*{Your turn}
\begin{itemize}
\raggedright
  \item \emph{Say it:} \emph{prôtibhā}
  \item \emph{Say it:} \emph{prôtibhā}, once more
  \item \emph{Say it:} say \emph{prôtibhā}
  \item \emph{Recall:} say the Bengali for knowledge, then the Bengali for courage, then say \emph{prôtibhā} again
\end{itemize}

\begin{tcolorbox}[breakable,colback=teal!4,colframe=teal!35!black,title={Before you move on}]
What does \bn{প্রতিভা} mean? (\textquotedblleft{}Talent\textquotedblright{}.) Say it once more.
\end{tcolorbox}
\section[\texorpdfstring{\bn{যন্ত্র} (jôntrô)}{যন্ত্র (jôntrô)}]{\bn{যন্ত্র} (jôntrô) — a machine, a tool}

\label{lesson:BN-C100-jontro}



\begin{quote}
Before the new one: say the Bengali for courage, then the Bengali for talent.
\end{quote}

\subsection*{You'll want to know: \bn{যন্ত্র}}
\textbf{\bn{যন্ত্র}} — \emph{jôntrô} — \textquotedblleft{}a machine, a tool\textquotedblright{}.

\begin{grammarlens}[title={What you've built}]
One more word for this chapter.
\end{grammarlens}

\subsection*{Your turn}
\begin{itemize}
\raggedright
  \item \emph{Say it:} \emph{jôntrô}
  \item \emph{Say it:} \emph{jôntrô}, once more
  \item \emph{Say it:} say \emph{jôntrô}
  \item \emph{Recall:} say the Bengali for courage, then the Bengali for talent, then say \emph{jôntrô} again
\end{itemize}

\begin{tcolorbox}[breakable,colback=teal!4,colframe=teal!35!black,title={Before you move on}]
What does \bn{যন্ত্র} mean? (\textquotedblleft{}A machine, a tool\textquotedblright{}.) Say it once more.
\end{tcolorbox}
\section[\texorpdfstring{\bn{সরঞ্জাম} (sôrônjām)}{সরঞ্জাম (sôrônjām)}]{\bn{সরঞ্জাম} (sôrônjām) — equipment}

\label{lesson:BN-C100-soronjam}



\begin{quote}
Before the new one: say the Bengali for talent, then the Bengali for a machine.
\end{quote}

\subsection*{You'll want to know: \bn{সরঞ্জাম}}
\textbf{\bn{সরঞ্জাম}} — \emph{sôrônjām} — \textquotedblleft{}equipment\textquotedblright{}.

\begin{grammarlens}[title={What you've built}]
One more word for this chapter.
\end{grammarlens}

\subsection*{Your turn}
\begin{itemize}
\raggedright
  \item \emph{Say it:} \emph{sôrônjām}
  \item \emph{Say it:} \emph{sôrônjām}, once more
  \item \emph{Say it:} say \emph{sôrônjām}
  \item \emph{Recall:} say the Bengali for talent, then the Bengali for a machine, then say \emph{sôrônjām} again
\end{itemize}

\begin{tcolorbox}[breakable,colback=teal!4,colframe=teal!35!black,title={Before you move on}]
What does \bn{সরঞ্জাম} mean? (\textquotedblleft{}Equipment\textquotedblright{}.) Say it once more.
\end{tcolorbox}
\section[\texorpdfstring{\bn{পদ্ধতি} (pôddhôti)}{পদ্ধতি (pôddhôti)}]{\bn{পদ্ধতি} (pôddhôti) — a method}

\label{lesson:BN-C100-poddhoti}



\begin{quote}
Before the new one: say the Bengali for a machine, then the Bengali for equipment.
\end{quote}

\subsection*{You'll want to know: \bn{পদ্ধতি}}
\textbf{\bn{পদ্ধতি}} — \emph{pôddhôti} — \textquotedblleft{}a method\textquotedblright{}.

\begin{grammarlens}[title={What you've built}]
One more word for this chapter.
\end{grammarlens}

\subsection*{Your turn}
\begin{itemize}
\raggedright
  \item \emph{Say it:} \emph{pôddhôti}
  \item \emph{Say it:} \emph{pôddhôti}, once more
  \item \emph{Say it:} say \emph{pôddhôti}
  \item \emph{Recall:} say the Bengali for a machine, then the Bengali for equipment, then say \emph{pôddhôti} again
\end{itemize}

\begin{tcolorbox}[breakable,colback=teal!4,colframe=teal!35!black,title={Before you move on}]
What does \bn{পদ্ধতি} mean? (\textquotedblleft{}A method\textquotedblright{}.) Say it once more.
\end{tcolorbox}
